\documentclass[../../main.tex]{subfiles}
\begin{document}

\begin{flushright}
\footnotesize\textit{【自動文字起こし・要確認】}
\end{flushright}

\noindent
[\textbf{解}] $n=2$ の時 $\left(\dfrac{3}{2}\right)^2 = \dfrac{9}{4} > 2 = 2!$ から成立. $n=k$ での成立を仮定する時
\[
f_k(x) = (k+1)\log \dfrac{k+1}{2} - \sum_{i=1}^k \log i
\]
とした時 $f_k(x) > 0$ である ($\because$ 与の同値).
\begin{align}
f_{k+1}(x) &= (k+2)\log \dfrac{k+2}{2} - \sum_{i=1}^{k+1} \log i \\
&= f_k(x) + \underbrace{(k+2)\log\dfrac{k+2}{2} - (k+1)\log\dfrac{k+1}{2} - \log(k+1)}_{g(k)} \quad \dots \text{①}
\end{align}

\noindent
ここで, $g(k) > 0$ を示す ($k \in \mathbb{R}$).
\begin{align}
g'(k) &= \log\dfrac{k+2}{2} + \dfrac{k+1}{k+2} - \left[ \log\dfrac{k+1}{2} + \dfrac{k}{k+1} \right] - \dfrac{1}{k+1} \\
&= \log\dfrac{k+2}{2} - \log\dfrac{k+1}{2} - \dfrac{1}{k+2} \\
g''(k) &= \dfrac{1}{k+2} - \dfrac{1}{k+1} + \dfrac{1}{(k+2)^2} \\
&= \dfrac{-1}{(k+1)(k+2)^2} < 0
\end{align}
より $g'(k)$ は単調減少. これと $g'(k) = \log\dfrac{k+2}{k+1} - \dfrac{1}{k+2} \to 0$ ($k \to \infty$) から $g'(k) > 0$ だから, $g(k)$ は単調増加であり, $k \ge 1$ の時
\[
g(1) = 2\log\dfrac{3}{2} - \log 2 = 2\log 3 - 3\log 2 = \log\dfrac{9}{8} > 0
\]
だから $g(k) > 0 \dots \text{②}$.
②と仮定から, ①より, $f_{k+1}(x) > 0$ となり $n=k+1$ でも成立.

\noindent
以上より示された.

\bigskip

\noindent
[\textbf{別 - もっと直接的に示す}]
\begin{align}
\left(\dfrac{k+2}{2}\right)^{k+1} - (k+1)! &= \dfrac{k+2}{2} \left(\dfrac{k+2}{k+1}\right)^k \left(\dfrac{k+1}{2}\right)^k - (k+1) \cdot k! \\
&\ge \left[ \dfrac{k+2}{2} \left(\dfrac{k+2}{k+1}\right)^k - (k+1) \right] k! \quad (\text{カテイ})
\end{align}
ここで, $k \ge 2$ のとき, 2項定理から
\[
\left(\dfrac{k+2}{k+1}\right)^k = \left(1 + \dfrac{1}{k+1}\right)^k \ge 1 + \dfrac{k}{k+1} = \dfrac{2k+1}{k+1}
\]
これから $> 0$ がわかるので OK!! \hfill $\mathbin{/\mkern-5mu/}$

\newpage

\end{document}
